\documentclass[12pt]{article}
\usepackage[margin=1in]{geometry}
\usepackage{setspace}
\onehalfspacing

\begin{document}

\noindent\today

\vspace{1cm}

\noindent Dear Editor,

\vspace{0.5cm}

I am pleased to submit the manuscript entitled ``Non-Directional Pre-Disclosure Drift Under Fair Disclosure: Evidence from Korean Earnings Announcements'' for consideration as an Original Article in the \textit{International Review of Financial Analysis}.

This study examines whether pre-disclosure price drift around Korean earnings announcements reflects information leakage. Using 2,221 DART filings from 64 KOSPI companies across 12 chaebol groups (2020--2026), we find that while statistically significant pre-disclosure drift exists (CAR[-5,$-$1] = +0.52\%, $p$<0.001), it carries no directional information about actual earnings surprises (49.0\% prediction accuracy, no better than chance). In contrast, the first post-disclosure trading day responds sharply and correctly (60.6\% accuracy, $p$<0.001).

The key contribution is distinguishing between the \textit{existence} of pre-disclosure drift and its \textit{directional informativeness}---a critical distinction for evaluating the effectiveness of electronic disclosure systems but rarely tested in the literature. We further provide cross-sectional analyses across 12 chaebol groups, firm size, and market beta that reveal significant heterogeneity but no pattern consistent with information leakage.

This work is relevant to IRFA readers interested in market microstructure, disclosure regulation, and emerging market efficiency. The manuscript has not been published or submitted elsewhere. All data are publicly available through DART, and analysis code is provided for full reproducibility.

\vspace{0.5cm}

\noindent Sincerely,

\vspace{0.5cm}

\noindent Yongjun Kim\\
Department of Software, Ajou University\\
yjkim0123@ajou.ac.kr

\end{document}
